\scanpage{45}
Calculons la \uline{trace} de la représentation \unsure{dérivée}, on trouve pour $\varphi\deleted{\unsure{\xi}}\in D/I$
\[
\operatorname{Tr}_{E^{\unsure{\natural|\xi}}}(\varphi)\deleted{\unsure{\xi}}=\sum_i\operatorname{Tr}_{E_i^{\unsure{\natural|\xi}}}(\varphi)
\]
Or si \deleted{$I$} $U$ opère trivialt sur $E$, \deleted{qui est} alors un module sur $D/U$, on a
\[
E^{\unsure{\natural|\xi}}=\operatorname{Im}\deleted{\unsure{\xi}}\pi_{\deleted{\unsure{I}}}\inserted{\text{projection}}\qquad
\pi=\frac{1}{N}\sum_{i\in I/U}e_i.\quad N=\operatorname{card} I/U
\]
donc si $g\in\unsure{I|D}/U\longrightarrow\varphi$ on a
\[
\operatorname{Tr}_{E^{\unsure{\natural|\xi}}}(\varphi)=\operatorname{Tr}_E(\deleted{\unsure{I}}\pi g)\deleted{=\frac{1}{N}\sum_{\substack{g\to\varphi\\\unsure{g\in I}}}\operatorname{Tr}_E(}
\]
d'où
\[
\boxed{\operatorname{Tr}_{E^{\unsure{\natural|\xi}}}(\varphi)=\frac{1}{N}\sum_{\substack{\psi_\alpha\in D/U\\\deleted{\unsure{g}}\\\psi_\alpha\to\varphi\in D/I}}\operatorname{Tr}_E\psi_\alpha}
\]
Comme \inserted{les deux mb. de} cette formule sont linéaires en $E$, ceci reste valable pour \unsure{tt} $E$ ; il faut simplement noter que $\operatorname{Tr}_E\psi$ est défini \deleted{\unsure{sur}} pour $\psi\in D/U$, si $U$ opère de façon unipotente, comme $\operatorname{Tr}_E(\psi')$ pour $\psi'\in D$ \unsure{se relevant} en $\psi$ : cette valeur ne \uline{dépend pas} du

\margin{44}

